\section{Related Work}
\label{sec:related_work}
Most ASIC- and FPGA-based LLM accelerators have focused on quantization or pruning to reduce the size of the model or memory usage.
\cite{hpca-2020-a3, isca-2021-elsa, hpca-2021-spatten, micro-2021-sanger, isca-2023-fact, asplos-2022-dota, isca-2022-mokey, isca-2022-leopard, micro-2022-butterfly, vlsi-2022-nmsparse} focused on the presence of unnecessary values in the attention matrix, which is the result of multiplying the Q and K matrices within the MHA layer of the transformer model. 
They proposed accelerators that approximate attention operations or prune unnecessary values in the attention matrix through methods such as eager prediction.
\cite{micro-2020-gobo, hpca-2021-spatten, isca-2023-olive} proposed accelerators that support quantization of the weight or embedding vector of the transformer model.
Meanwhile, \cite{micro-2022-dfx, asplos-2023-flat} proposed accelerators that process the transformer model losslessly.
\cite{micro-2022-dfx} proposed a system that distributes processing across multiple FPGAs for low-latency inference of the transformer model and \cite{asplos-2023-flat} proposed a dataflow for the accelerator to improve the efficiency of the attention operation.

Numerous studies have proposed leveraging near-memory processing or PIM to process the transformer model effectively.
\cite{isca-2021-hbm-pim, isscc-2022-aim} introduced commercialized PIM architectures on HBM and GDDR, respectively.
They proposed a system and software stack that offloads the GEMV operations of the transformer model to the PIM architecture with in-bank processing units.
\cite{asplos-2024-attacc} proposed a PIM architecture that efficiently handles attention operations with extremely low-\opb and high memory capacity requirements on large batch sizes.
\cite{asplos-2024-neupim} proposed the integration of NPU and PIMs to efficiently handle both GEMM and GEMV.
\cite{hpca-2022-transpim} proposed an end-to-end accelerator to process the transformer model on PIM.

TOP-PIM~\cite{hpdc-2014-toppim} proposed GPU with PIM stacks, achieving acceleration for memory-intensive workloads through the internal stack bandwidth between the DRAM die and logic die, which is higher than the I/O bandwidth of the HBM stack.
However, it did not aim to accelerate LLMs and lacks details of DRAM die architectures.
Also, it did not consider concurrently executing GPUs and PIM stacks; thus, the expert and attention co-processing cannot be applied.
%
Despite these works on accelerating LLMs and PIM architectures, to the best of our knowledge, \duplex is the first to accelerate LLMs with MoE and GQA in continuous batching.


\label{sec:related_work:accelerator}

